% Contenu éditorial du cas Gryffondor--Serpentard.
% Toutes les valeurs numériques affichées sont fournies par le générateur.

% Maisons et matières portent partout la teinte qui les désigne dans les
% figures. Les noms sont écrits en entier chaque fois que la place le permet ;
% les tableaux les plus serrés emploient les formes abrégées de même couleur.
\newcommand{\GryffondorName}{\textcolor{coursescore!88!ink}{\MicroGryffondorLabel}}
\newcommand{\SerpentardName}{\textcolor{warning!88!ink}{\MicroSerpentardLabel}}
\newcommand{\PotionsName}{\textcolor{coursepotions!92!ink}{Potions}}
\newcommand{\VolName}{\textcolor{coursevol!92!ink}{Vol}}
\newcommand{\microheader}{\rowcolor{ink!92}\color{white}\bfseries}

% Les deux notes vivent sur des échelles différentes : sans « scale only axis »,
% la boîte s'étire et les deux groupes d'élèves se lisent comme des traînées
% obliques. Le rapport 1,25 entre largeur et hauteur rend un disque de notes
% rond à l'écran.
\pgfplotsset{
  rawcloud/.style={
    scale only axis,width=7cm,height=5.6cm,
    xmin=0,xmax=20,ymin=0,ymax=100,
    xtick={0,5,10,15,20},ytick={0,20,40,60,80,100},
    axis lines=box,axis line style={ink!42},
    grid=major,grid style={gridgray!34},
    xlabel={\PotionsName{} (sur 20)},ylabel={\VolName{} (sur 100)},
    tick label style={font=\scriptsize},label style={font=\scriptsize},
    legend style={at={(.5,1.02)},anchor=south,draw=none,
      font=\scriptsize,legend columns=3,
      /tikz/every even column/.append style={column sep=7pt}}},
  smallcloud/.style={
    rawcloud,width=4.5cm,height=3.6cm,
    xtick={0,10,20},ytick={0,50,100},
    xlabel={\PotionsName},ylabel={\VolName},
    tick label style={font=\tiny}},
  stdcloud/.style={
    scale only axis,width=4.5cm,height=3.6cm,
    xmin=-2.6,xmax=2.6,ymin=-2.6,ymax=2.6,
    xtick={-2,-1,0,1,2},ytick={-2,-1,0,1,2},
    axis lines=box,axis line style={ink!42},
    grid=major,grid style={gridgray!34},
    xlabel={$x_{\mathrm{pot}}$},ylabel={$x_{\mathrm{vol}}$},
    tick label style={font=\tiny},label style={font=\scriptsize}}}

% Vignettes de la planche d'ensemble. Chaque carte est une image autonome, posée
% par un nœud : son contenu est donc centré par construction, jamais par des
% coordonnées ajustées à la main. L'ensemble est enveloppé dans \hyperref, ce qui
% rend toute la surface de la carte cliquable vers la section correspondante.
\newlength{\chainw}\setlength{\chainw}{3.06cm}
\newlength{\chainh}\setlength{\chainh}{3.6cm}
\pgfplotsset{
  chainaxis/.style={
    scale only axis,width=2.55cm,height=1.7cm,
    axis lines=box,axis line style={ink!30,line width=.4pt},
    xtick=\empty,ytick=\empty,
    clip mode=individual,
    every axis plot/.append style={line width=.7pt}}}

% Lignes de niveau de la perte dans le plan des deux poids, calculées par le
% générateur. Elles servent de fond commun aux vignettes gradient et mise à
% jour : l'une y lit une direction, l'autre y suit une suite de pas.
\newcommand{\microrelief}{%
  \addplot[courseupdate!22!white,line width=.5pt] coordinates {\MicroContourOuter};
  \addplot[courseupdate!34!white,line width=.5pt] coordinates {\MicroContourThird};
  \addplot[courseupdate!48!white,line width=.5pt] coordinates {\MicroContourSecond};
  \addplot[courseupdate!62!white,line width=.5pt] coordinates {\MicroContourInner};}

\newcommand{\chaincard}[4]{%
  \hyperref[#1]{%
    \begin{tikzpicture}
      \path[draw=ink!24,fill=white,rounded corners=1.6pt,line width=.5pt]
        (0,0) rectangle (\chainw,\chainh);
      \path[fill=#2!16,rounded corners=1.6pt]
        (0,\chainh-.51cm) rectangle (\chainw,\chainh);
      \node[font=\fontsize{6.3}{7}\selectfont\bfseries,text=ink!92,anchor=west]
        at (.16,\chainh-.255cm) {#3};
      \node[inner sep=0pt] at (.5\chainw,.5\chainh-.255cm) {#4};
    \end{tikzpicture}}}

% Frontière ajustée et frontière décalée par le seuil haut, toutes deux issues
% des coefficients estimés : aucune constante n'est écrite ici à la main.
\newcommand{\microboundary}[1]{%
  \addplot[#1,domain=0:20,samples=2]
    {\MicroBoundarySlope*x+(\MicroBoundaryIntercept)}}
\newcommand{\microboundaryhigh}[1]{%
  \addplot[#1,domain=0:20,samples=2]
    {\MicroBoundarySlope*x+(\MicroBoundaryInterceptHigh)}}

\courseintro

\begin{frame}[t]{Formulation du problème de classification}
  \vspace{4mm}
  \begin{columns}[T,onlytextwidth]
    \column{.47\textwidth}
      \begin{coursecallout}{coursedata!88!ink}{Observations disponibles}
        Pour chaque élève :
        \begin{itemize}
          \item une note de \PotionsName{} sur 20 ;
          \item une note de \VolName{} sur 100.
        \end{itemize}
      \end{coursecallout}
    \column{.47\textwidth}
      \begin{coursecallout}{coursedecision!88!ink}{Maison observée}
        Deux maisons sont représentées :
        \begin{itemize}
          \item \GryffondorName ;
          \item \SerpentardName.
        \end{itemize}
      \end{coursecallout}
  \end{columns}
  \vfill
  \centering
  {\large\bfseries Construire une règle qui associe les deux notes à une maison.}
\end{frame}

% Les titres de planche passent par \MakeUppercase : aucune commande colorée ne
% peut y figurer, le nom de la couleur serait mis en capitales lui aussi.
\begin{frame}[t]{Distribution des notes de Potions et de Vol}
  \begin{columns}[c,onlytextwidth]
    \column{.34\textwidth}
      \scriptsize
      \renewcommand{\arraystretch}{.66}
      \setlength{\tabcolsep}{3pt}
      \begin{tabular}{@{}l l r r@{}}
        \toprule
        élève & maison & \PotionsName & \VolName\\
        \midrule
        \MicroCompactRows
        \bottomrule
      \end{tabular}
    \column{.64\textwidth}
      \centering
      \begin{tikzpicture}
        \begin{axis}[rawcloud,width=7.15cm,height=5.72cm]
          \addplot[only marks,mark=square*,mark size=2.6pt,coursescore!88!ink]
            coordinates {\MicroAllGryffondorCoords};
          \addlegendentry{\GryffondorName}
          \addplot[only marks,mark=*,mark size=2.6pt,warning!88!ink]
            coordinates {\MicroAllSerpentardCoords};
          \addlegendentry{\SerpentardName}
          \addplot[only marks,mark=o,mark size=4.6pt,very thick,ink!72]
            coordinates {\MicroMissingCoords};
          \addlegendentry{note imputée}
          \MicroAllNameLabels
        \end{axis}
      \end{tikzpicture}
  \end{columns}
\end{frame}

\begin{frame}[t]{Chaîne de construction du classifieur}
\vfill
\centering
\begin{adjustbox}{max totalsize={\textwidth}{.83\textheight},center}
\begin{tikzpicture}[
  flow/.style={-{Latex[length=1.8mm]},line width=.75pt,ink!62},
  slot/.style={inner sep=0pt,outer sep=0pt}]

  \node[slot] (donnees) at (0,4.8) {\chaincard{course-stage-start-1}{coursedata}{DONNÉES}{%
    \begin{tikzpicture}
      \begin{axis}[chainaxis,xmin=0,xmax=20,ymin=0,ymax=100,
        xlabel={\PotionsName},ylabel={\VolName},
        label style={font=\tiny,inner sep=1.5pt}]
        \addplot[only marks,mark=square*,mark size=.9pt,coursescore!88!ink]
          coordinates {\MicroAllGryffondorCoords};
        \addplot[only marks,mark=*,mark size=.9pt,warning!88!ink]
          coordinates {\MicroAllSerpentardCoords};
      \end{axis}
    \end{tikzpicture}}};

  \node[slot] (score) at (3.78,4.8) {\chaincard{course-stage-start-2}{coursescore}{SCORE LINÉAIRE}{%
    \begin{tikzpicture}
      \begin{axis}[chainaxis,xmin=0,xmax=20,ymin=0,ymax=100]
        \addplot[only marks,mark=square*,mark size=.9pt,coursescore!88!ink]
          coordinates {\MicroTrainGryffondorCoords};
        \addplot[only marks,mark=*,mark size=.9pt,warning!88!ink]
          coordinates {\MicroTrainSerpentardCoords};
        \microboundary{ink!72,line width=.8pt};
      \end{axis}
    \end{tikzpicture}}};

  \node[slot] (sigmoide) at (7.56,4.8) {\chaincard{course-stage-start-3}{courseprobability}{SIGMOÏDE}{%
    \begin{tikzpicture}
      \begin{axis}[chainaxis,xmin=-3.2,xmax=3.2,ymin=-.06,ymax=1.06]
        \addplot[courseprobability!88!ink,domain=-3.2:3.2,samples=70]
          {1/(1+exp(-x))};
        \addplot[only marks,mark=square*,mark size=.9pt,coursescore!88!ink]
          coordinates {\MicroSigmoidGryffondorCoords};
        \addplot[only marks,mark=*,mark size=.9pt,warning!88!ink]
          coordinates {\MicroSigmoidSerpentardCoords};
      \end{axis}
    \end{tikzpicture}}};

  \node[slot] (decision) at (11.34,4.8) {\chaincard{course-stage-start-7}{coursedecision}{DÉCISION}{%
    \begin{tikzpicture}
      \begin{axis}[chainaxis,height=.95cm,xmin=0,xmax=1,ymin=0,ymax=1]
        \fill[warning!8] (axis cs:0,0) rectangle (axis cs:.5,1);
        \fill[coursescore!8] (axis cs:.5,0) rectangle (axis cs:1,1);
        \draw[coursedecision!88!ink,dashed,line width=.6pt]
          (axis cs:.5,0) -- (axis cs:.5,1);
        \addplot[only marks,mark=square*,mark size=.9pt,coursescore!88!ink]
          coordinates {\MicroDecisionGryffondorCoords};
        \addplot[only marks,mark=*,mark size=.9pt,warning!88!ink]
          coordinates {\MicroDecisionSerpentardCoords};
        \node[font=\tiny,text=warning!88!ink,anchor=north west,yshift=-1pt]
          at (axis description cs:0,0) {\MicroSerpentardLabel};
        \node[font=\tiny,text=coursescore!88!ink,anchor=north east,yshift=-1pt]
          at (axis description cs:1,0) {\MicroGryffondorLabel};
        \node[font=\tiny,text=coursedecision!88!ink,anchor=north,
          inner sep=1.5pt] at (axis description cs:.5,1) {$\tau$};
      \end{axis}
    \end{tikzpicture}}};

  \node[slot] (maj) at (0,0) {\chaincard{course-stage-start-6}{courseupdate}{MISE À JOUR}{%
    \begin{tikzpicture}
      \begin{axis}[chainaxis,xmin=\MicroSliceXmin,xmax=\MicroSliceXmax,
        ymin=\MicroSliceYmin,ymax=\MicroSliceYmax,
        xlabel={$w_{\mathrm{pot}}$},ylabel={$w_{\mathrm{vol}}$},
        label style={font=\tiny,inner sep=1.5pt}]
        \microrelief
        \addplot[courseupdate!85!ink,line width=.8pt,mark=*,mark size=.75pt,
          mark options={fill=courseupdate!85!ink,draw=none}]
          coordinates {\MicroGDPathCoords};
        \addplot[only marks,mark=star,mark size=2pt,line width=.7pt,ink!75]
          coordinates {\MicroOptimumCoords};
      \end{axis}
    \end{tikzpicture}}};

  \node[slot] (gradient) at (3.78,0) {\chaincard{course-stage-start-5}{coursegradient}{GRADIENT}{%
    \begin{tikzpicture}
      \begin{axis}[chainaxis,xmin=\MicroSliceXmin,xmax=\MicroSliceXmax,
        ymin=\MicroSliceYmin,ymax=\MicroSliceYmax,
        xlabel={$w_{\mathrm{pot}}$},ylabel={$w_{\mathrm{vol}}$},
        label style={font=\tiny,inner sep=1.5pt}]
        \microrelief
        \draw[-{Latex[length=1.7mm]},coursegradient!92!ink,line width=1.1pt]
          (axis cs:\MicroGradientProbeX,\MicroGradientProbeY)
          -- (axis cs:\MicroGradientTipX,\MicroGradientTipY);
        \addplot[only marks,mark=*,mark size=1.2pt,coursegradient!92!ink]
          coordinates {\MicroGradientProbe};
        \node[font=\tiny,text=coursegradient!92!ink,anchor=south west,inner sep=1pt]
          at (axis cs:\MicroGradientTipX,\MicroGradientTipY) {$-\nabla J$};
      \end{axis}
    \end{tikzpicture}}};

  \node[slot] (perte) at (7.56,0) {\chaincard{course-stage-start-4}{courseloss}{PERTE LOGISTIQUE}{%
    \begin{tikzpicture}
      \begin{axis}[chainaxis,xmin=-3,xmax=3,ymin=0,ymax=3.2,
        xlabel={$z$},label style={font=\tiny,inner sep=1.5pt}]
        \addplot[coursescore!88!ink,domain=-3:3,samples=60] {ln(1+exp(-x))};
        \addplot[warning!88!ink,domain=-3:3,samples=60] {ln(1+exp(x))};
        \node[font=\fontsize{5}{5.5}\selectfont,text=coursescore!88!ink,
          anchor=north west,inner sep=2pt]
          at (axis description cs:0,1) {\MicroGryffondorLabel};
        \node[font=\fontsize{5}{5.5}\selectfont,text=warning!88!ink,
          anchor=north east,inner sep=2pt]
          at (axis description cs:1,1) {\MicroSerpentardLabel};
      \end{axis}
    \end{tikzpicture}}};

  \node[slot] (evaluation) at (11.34,0)
    {\chaincard{course-stage-start-8}{courseevaluation}{ÉVALUATION}{%
    \begin{tikzpicture}[x=1.02cm,y=.58cm,
      cell/.style={draw=white,line width=.7pt,minimum width=1cm,
        minimum height=.56cm,font=\tiny\bfseries,inner sep=0pt},
      tag/.style={font=\tiny,text=ink!62}]
      \node[cell,fill=rulecolor!18,text=rulecolor!75!ink] at (0,1) {VN $=\MicroTrainTN$};
      \node[cell,fill=warning!16,text=warning!85!ink]     at (1,1) {FP $=\MicroTrainFP$};
      \node[cell,fill=warning!16,text=warning!85!ink]     at (0,0) {FN $=\MicroTrainFN$};
      \node[cell,fill=rulecolor!18,text=rulecolor!75!ink] at (1,0) {VP $=\MicroTrainTP$};
      \node[tag,anchor=south] at (.5,1.62) {maison prédite};
      \node[tag,anchor=south,rotate=90] at (-.62,.5) {maison observée};
    \end{tikzpicture}}};

  \draw[flow] (donnees) -- (score);
  \draw[flow] (score) -- (sigmoide);
  \draw[flow] (sigmoide) -- (decision);
  \draw[flow] (sigmoide) -- (perte);
  \draw[flow] (perte) -- (gradient);
  \draw[flow] (gradient) -- (maj);
  \draw[flow] (decision) -- (evaluation);
  \draw[flow] (maj.north) -- ++(0,.55) -| (score.south);
\end{tikzpicture}
\end{adjustbox}
\vfill
\end{frame}

% ---------------------------------------------------------------------------
\coursepart{Définir les données et le protocole}{Données}{coursedata}

\begin{frame}[t]{Séparation entre estimation et évaluation}
  \vspace{5mm}
  \centering
  \begin{tikzpicture}[node distance=9mm and 12mm,
    box/.style={draw=ink!28,rounded corners=2pt,minimum height=15mm,
      minimum width=31mm,align=center,font=\small,inner xsep=4mm},
    arr/.style={-{Latex[length=2mm]},thick,ink!58}]
    \node[box,fill=coursedata!13] (all) {24 élèves};
    \node[box,fill=coursedata!22,right=of all,yshift=10mm] (train)
      {\MicroNTrain{} élèves\\estimation};
    \node[box,fill=courseevaluation!18,right=of all,yshift=-10mm] (test)
      {\MicroNTest{} élèves\\évaluation};
    \node[box,fill=courseupdate!15,right=of train] (fit)
      {médianes, $\mu$, $s$\\et coefficients};
    \node[box,fill=courseevaluation!24,right=of test] (eval)
      {mesure\\des erreurs};
    \draw[arr] (all) -- (train);
    \draw[arr] (all) -- (test);
    \draw[arr] (train) -- (fit);
    \draw[arr] (fit.south) |- (eval.north);
    \draw[arr] (test) -- (eval);
  \end{tikzpicture}
  \vfill
  Le groupe d'estimation constitue le \emph{jeu d'apprentissage} : lui seul fixe
  les nombres du modèle. Le groupe d'évaluation les reçoit sans les modifier.
\end{frame}

\begin{frame}[t]{Groupe d'estimation : \MicroNTrain{} élèves}
  \centering
  \tiny
  \renewcommand{\arraystretch}{.9}
  \begin{tabularx}{\textwidth}{c >{\raggedright\arraybackslash}X l c c}
    \microheader ID & \color{white}\bfseries Élève &
      \color{white}\bfseries Maison & \color{white}\bfseries Potions &
      \color{white}\bfseries Vol \\
    \MicroTrainRows
  \end{tabularx}
  \par\vspace{2mm}
  \begin{vigilance}[Une note manque ici même]
    La note de \PotionsName{} de Ron est absente. Le trou se trouve donc là où sont
    calculées les statistiques : sa valeur de remplacement participera aux
    médianes, aux moyennes et aux coefficients.
  \end{vigilance}
\end{frame}

\begin{frame}[t]{Groupe d'évaluation : huit élèves}
  \centering
  \small
  \renewcommand{\arraystretch}{1.04}
  \begin{tabularx}{\textwidth}{c >{\raggedright\arraybackslash}X l c c}
    \microheader ID & \color{white}\bfseries Élève &
      \color{white}\bfseries Maison & \color{white}\bfseries Potions &
      \color{white}\bfseries Vol \\
    \MicroTestRawRows
  \end{tabularx}
  \vfill
  \begin{vigilance}[Une note manque aussi à l'évaluation]
    La note de \VolName{} de George est absente. Une case vide ne peut pas être
    standardisée : elle doit d'abord recevoir une valeur numérique, et cette
    valeur vient du groupe d'estimation.
  \end{vigilance}
\end{frame}

\begin{frame}[t]{Compléter, puis standardiser une note manquante}
  \vspace{3mm}
  \begin{columns}[T,onlytextwidth]
    \column{.47\textwidth}
      \begin{coursecallout}{coursedata!88!ink}{Valeurs de remplacement}
        \[
          \widetilde P=\MicroMedianPotions,
          \qquad \widetilde V=\MicroMedianVol.
        \]
        Ces médianes ne sont lues que sur les notes \emph{observées} du groupe
        d'estimation, jamais sur le groupe d'évaluation.
      \end{coursecallout}
    \column{.47\textwidth}
      \begin{coursecallout}{courseevaluation!88!ink}{Deux substitutions}
        \[
          \begin{aligned}
            P_{\text{Ron}}&\leftarrow\MicroMedianPotions
              &\Longrightarrow x_{\mathrm{pot}}&=0,\\
            V_{\text{George}}&\leftarrow\MicroMedianVol
              &\Longrightarrow x_{\mathrm{vol}}&=0{,}35.
          \end{aligned}
        \]
        En \PotionsName{} la médiane égale la moyenne ; en \VolName{} elle vaut
        \(\MicroMedianVol\) contre une moyenne de \(\MicroMeanVol\).
      \end{coursecallout}
  \end{columns}
  \vfill
  \centering
  \begin{adjustbox}{max width=.96\textwidth,center}
  \begin{tikzpicture}[node distance=7mm,
    box/.style={draw=ink!25,rounded corners=2pt,minimum height=11mm,
      align=center,font=\small,inner xsep=4mm},
    arr/.style={-{Latex[length=1.8mm]},thick,ink!58}]
    \node[box,fill=warning!12] (raw) {note manquante};
    \node[box,fill=coursedata!16,right=of raw] (rule) {médiane\\de la matière};
    \node[box,fill=rulecolor!16,right=of rule] (ready) {valeur\\numérique $v$};
    \node[box,fill=courseprobability!14,right=of ready] (standard)
      {$x=(v-\mu)/s$\\note standardisée};
    \draw[arr] (raw) -- (rule);
    \draw[arr] (rule) -- (ready);
    \draw[arr] (ready) -- (standard);
  \end{tikzpicture}
  \end{adjustbox}
\end{frame}

\begin{frame}[t]{Deux matières, deux échelles de mesure}
  \vspace{2mm}
  \begin{columns}[T,onlytextwidth]
    \column{.46\textwidth}
      \centering
      {\bfseries\PotionsName}
      \begin{tikzpicture}[y=.045cm]
        \draw[-{Latex[length=1.7mm]},ink!58] (0,0) -- (0,23);
        \foreach \v in {0,5,10,15,20}
          \draw[ink!55] (-.12,\v) -- (.12,\v)
            node[right=2mm,font=\scriptsize] {\v};
        \fill[coursepotions!62] (-.35,\MicroPotionsMin) rectangle (.35,\MicroPotionsMax);
      \end{tikzpicture}
      \par notée sur 20
    \column{.46\textwidth}
      \centering
      {\bfseries\VolName}
      \begin{tikzpicture}[y=.009cm]
        \draw[-{Latex[length=1.7mm]},ink!58] (0,0) -- (0,113);
        \foreach \v in {0,25,50,75,100}
          \draw[ink!55] (-.12,\v) -- (.12,\v)
            node[right=2mm,font=\scriptsize] {\v};
        \fill[coursevol!55] (-.35,\MicroVolMin) rectangle (.35,\MicroVolMax);
      \end{tikzpicture}
      \par notée sur 100
  \end{columns}
  \vfill
  Un point représente 5\% de l'échelle en \PotionsName, contre 1\% en \VolName.
\end{frame}

\begin{frame}[t]{Centrage-réduction des deux notes}
  \vspace{2mm}
  Les paramètres de standardisation sont calculés sur le groupe d'estimation,
  puis conservés :
  \begin{columns}[T,onlytextwidth]
    \column{.47\textwidth}
      \begin{coursecallout}{coursepotions!92!ink}{\PotionsName}
        \[
          \mu_{\mathrm{pot}}=\MicroMeanPotions,
          \quad s_{\mathrm{pot}}=\MicroSdPotions,
        \]
        \[
          x_{\mathrm{pot}}=
          \frac{P-\MicroMeanPotions}{\MicroSdPotions}.
        \]
      \end{coursecallout}
    \column{.47\textwidth}
      \begin{coursecallout}{coursevol!92!ink}{\VolName}
        \[
          \mu_{\mathrm{vol}}=\MicroMeanVol,
          \quad s_{\mathrm{vol}}=\MicroSdVol,
        \]
        \[
          x_{\mathrm{vol}}=
          \frac{V-\MicroMeanVol}{\MicroSdVol}.
        \]
      \end{coursecallout}
  \end{columns}
  \vfill
  Pour Harry : \(x_{\mathrm{pot}}=\MicroHarryPotions\) et
  \(x_{\mathrm{vol}}=\MicroHarryVol\). Les deux coordonnées sont désormais sans
  unité.
  \par\smallskip
  La covariance d'estimation vaut \(\MicroTrainCovariance\), soit une
  corrélation de \(\MicroTrainCorrelation\) : les deux notes varient en sens
  opposé.
\end{frame}

\begin{frame}[t]{Effet géométrique de la standardisation}
  \centering
  \begin{columns}[T,onlytextwidth]
    \column{.49\textwidth}
      \centering\small Notes brutes
      \begin{tikzpicture}
        \begin{axis}[smallcloud]
          \addplot[only marks,mark=square*,mark size=2.2pt,coursescore]
            coordinates {\MicroTrainGryffondorCoords};
          \addplot[only marks,mark=*,mark size=2.2pt,warning]
            coordinates {\MicroTrainSerpentardCoords};
        \end{axis}
      \end{tikzpicture}
    \column{.49\textwidth}
      \centering\small Coordonnées standardisées
      \begin{tikzpicture}
        \begin{axis}[stdcloud]
          \addplot[only marks,mark=square*,mark size=2.2pt,coursescore]
            coordinates {\MicroTrainGryffondorStdCoords};
          \addplot[only marks,mark=*,mark size=2.2pt,warning]
            coordinates {\MicroTrainSerpentardStdCoords};
        \end{axis}
      \end{tikzpicture}
  \end{columns}
  \vfill
  Le centrage-réduction change les unités, pas l'identité des élèves ni leur maison.
\end{frame}

% ---------------------------------------------------------------------------
\coursepart{Construire le score linéaire}{Score linéaire}{coursescore}

\begin{frame}[t]{Définition du score linéaire}
  \vspace{5mm}
  \[
    \boxed{z_i=w_0+w_{\mathrm{pot}}x_{i,\mathrm{pot}}
                    +w_{\mathrm{vol}}x_{i,\mathrm{vol}}}
  \]
  \vspace{4mm}
  \begin{columns}[T,onlytextwidth]
    \column{.47\textwidth}
      \begin{coursecallout}{coursescore!88!ink}{Coefficients}
        Le signe indique le sens de l'association ; la valeur absolue en mesure
        l'intensité sur l'échelle standardisée.
      \end{coursecallout}
    \column{.47\textwidth}
      \begin{coursecallout}{coursegradient!88!ink}{État du modèle}
        Les coefficients sont encore inconnus. Leur estimation viendra après
        la définition de la perte et du gradient.
      \end{coursecallout}
  \end{columns}
\end{frame}

\begin{frame}[t]{Contributions des deux matières au score}
  \centering
  Prenons provisoirement \(w^{\mathrm{essai}}=\MicroTrialWeights\), donc
  \(\MicroTrialScore\).
  \vspace{3mm}
  \scriptsize
  \renewcommand{\arraystretch}{1.22}
  \begin{tabular}{lrrrrrrr}
    \toprule
    élève & $x_{\mathrm{pot}}$ & $x_{\mathrm{vol}}$ &
      $-2x_{\mathrm{pot}}$ & $+x_{\mathrm{vol}}$ & $z$ & $p$ & $\ell$\\
    \midrule
    \MicroTrialRows
    \bottomrule
  \end{tabular}
  \vfill
  \normalsize
  Le score agrège les contributions ; il ne constitue encore ni une
  probabilité ni une classe prédite.
\end{frame}

\begin{frame}[t]{Le niveau de score nul définit une droite}
  \begin{columns}[c,onlytextwidth]
    \column{.60\textwidth}
      \begin{tikzpicture}
        \begin{axis}[rawcloud,legend columns=2]
          \addplot[only marks,mark=square*,mark size=2.6pt,coursescore!88!ink]
            coordinates {\MicroTrainGryffondorCoords};
          \addlegendentry{\GryffondorName}
          \addplot[only marks,mark=*,mark size=2.6pt,warning!88!ink]
            coordinates {\MicroTrainSerpentardCoords};
          \addlegendentry{\SerpentardName}
          \microboundary{ink!78,very thick};
          \node[font=\scriptsize,text=coursescore!88!ink,anchor=south]
            at (axis cs:4,74) {$z>0$};
          \node[font=\scriptsize,text=warning!88!ink,anchor=north]
            at (axis cs:16,24) {$z<0$};
        \end{axis}
      \end{tikzpicture}
    \column{.36\textwidth}
      \begin{coursecallout}{coursescore!88!ink}{Une droite, deux demi-plans}
        L'ensemble \(z=0\) est une droite du plan des notes ; elle le partage
        en deux demi-plans où le score est de signe constant.
      \end{coursecallout}
      \par\smallskip
      {\small Après estimation des coefficients, son équation sera
      \(\MicroBoundaryRaw\).}
  \end{columns}
\end{frame}

\begin{frame}[t]{Écriture matricielle des scores}
  Pour les \MicroNTrain{} observations du groupe d'estimation :
  \[
    X=
    \begin{pmatrix}
      1 & x_{1,\mathrm{pot}} & x_{1,\mathrm{vol}}\\
      \vdots & \vdots & \vdots\\
      1 & x_{\MicroNTrain,\mathrm{pot}} & x_{\MicroNTrain,\mathrm{vol}}
    \end{pmatrix},
    \qquad
    w=\begin{pmatrix}w_0\\w_{\mathrm{pot}}\\w_{\mathrm{vol}}\end{pmatrix}.
  \]
  \[
    \boxed{z=Xw}
  \]
  \vfill
  Une même multiplication produit les \MicroNTrain{} scores ; chaque ligne de
  \(X\) correspond à un élève.
\end{frame}

% ---------------------------------------------------------------------------
\coursepart{Transformer le score en probabilité}{Sigmoïde}{courseprobability}

\begin{frame}[t]{La sigmoïde associe une probabilité à tout score réel}
  \centering
  \begin{tikzpicture}
    \begin{axis}[
      width=.78\textwidth,height=5.5cm,
      xmin=-5,xmax=5,ymin=-.05,ymax=1.05,
      axis lines=middle,grid=major,grid style={gridgray!35},
      xlabel={$z$},ylabel={$p$},
      xtick={-4,-2,0,2,4},ytick={0,.5,1},
      tick label style={font=\scriptsize},label style={font=\scriptsize}]
      \addplot[courseprobability!88!ink,very thick,domain=-5:5,samples=120]
        {1/(1+exp(-x))};
      \addplot[ink!45,dashed,domain=-5:5,samples=2] {.5};
      \addplot[only marks,mark=*,mark size=2.5pt,coursescore]
        coordinates {(0,.5)};
    \end{axis}
  \end{tikzpicture}
  \[
    \boxed{p_i=\sigma(z_i)=\frac{1}{1+e^{-z_i}}}
  \]
\end{frame}

\begin{frame}[t]{Propriétés de la fonction logistique}
  \vspace{4mm}
  \begin{columns}[T,onlytextwidth]
    \column{.31\textwidth}
      \begin{coursecallout}{courseprobability!88!ink}{Monotonie}
        \[
          z_1<z_2\Rightarrow\sigma(z_1)<\sigma(z_2).
        \]
      \end{coursecallout}
    \column{.31\textwidth}
      \begin{coursecallout}{coursescore!88!ink}{Point central}
        \[
          \sigma(0)=\frac12.
        \]
      \end{coursecallout}
    \column{.31\textwidth}
      \begin{coursecallout}{courseloss!88!ink}{Symétrie}
        \[
          \sigma(-z)=1-\sigma(z).
        \]
      \end{coursecallout}
  \end{columns}
  \vfill
  \centering
  La sigmoïde conserve l'ordre des scores et les transforme en valeurs
  strictement comprises entre \(0\) et \(1\).
\end{frame}

\begin{frame}[t]{Le logit est l'inverse de la sigmoïde}
  \vspace{5mm}
  \begin{columns}[T,onlytextwidth]
    \column{.47\textwidth}
      \begin{coursecallout}{courseprobability!88!ink}{Du score à la probabilité}
        \[
          \frac{p}{1-p}=e^z,
          \qquad p=\frac{e^z}{1+e^z}.
        \]
      \end{coursecallout}
    \column{.47\textwidth}
      \begin{coursecallout}{coursescore!88!ink}{De la probabilité au score}
        \[
          \operatorname{logit}(p)
          =\ln\!\left(\frac{p}{1-p}\right)=z.
        \]
      \end{coursecallout}
  \end{columns}
  \vfill
  \centering
  \begin{tabular}{c|ccccc}
    $z$ & $-2$ & $-1$ & $0$ & $1$ & $2$\\
    \midrule
    $p$ & $0{,}119$ & $0{,}269$ & $1/2$ & $0{,}731$ & $0{,}881$
  \end{tabular}
\end{frame}

\begin{frame}[t]{Probabilités associées aux scores d'essai}
  \centering
  \renewcommand{\arraystretch}{1.35}
  \begin{tabular}{llccl}
    \toprule
    élève & maison observée & $z$ d'essai & $p=\sigma(z)$ & maison favorisée\\
    \midrule
    \MicroTrialProbabilityRows
    \bottomrule
  \end{tabular}
  \vfill
  \begin{coursecallout}{courseprobability!88!ink}{}
    La sigmoïde étant strictement croissante, deux observations de même score
    reçoivent la même probabilité selon le modèle. Les lignes marquées
    \textcolor{warning}{$\bullet$} sont celles pour lesquelles la maison
    favorisée contredit la maison observée.
  \end{coursecallout}
\end{frame}

% ---------------------------------------------------------------------------
\coursepart{Définir la perte logistique}{Perte logistique}{courseloss}

\begin{frame}[t]{Codage de la maison observée}
  \vspace{5mm}
  \[
    y_i=
    \begin{cases}
      1 & \text{si l'élève appartient à \GryffondorName},\\
      0 & \text{si l'élève appartient à \SerpentardName}.
    \end{cases}
  \]
  \vspace{5mm}
  \begin{resultat}
    Ce codage n'était nécessaire ni pour calculer le score \(z_i\), ni pour le
    transformer en probabilité \(p_i\). Il devient nécessaire pour comparer la
    sortie du modèle à la maison observée.
  \end{resultat}
\end{frame}

\begin{frame}[t]{Perte logistique d'une observation}
  \vspace{3mm}
  \[
    \boxed{\ell(z,y)=-y\ln p-(1-y)\ln(1-p)}
  \]
  \begin{columns}[T,onlytextwidth]
    \column{.47\textwidth}
      \begin{coursecallout}{coursescore!88!ink}{Maison observée : \GryffondorName}
        \[
          y=1,\qquad \ell=-\ln p=\ln(1+e^{-z}).
        \]
        La perte décroît lorsque le score augmente.
      \end{coursecallout}
    \column{.47\textwidth}
      \begin{coursecallout}{warning!88!ink}{Maison observée : \SerpentardName}
        \[
          y=0,\qquad \ell=-\ln(1-p)=\ln(1+e^z).
        \]
        La perte décroît lorsque le score diminue.
      \end{coursecallout}
  \end{columns}
\end{frame}

\begin{frame}[t]{Comportement de la perte selon la maison observée}
  \centering
  \begin{tikzpicture}
    \begin{axis}[width=.74\textwidth,height=5.2cm,xmin=-4,xmax=4,ymin=0,ymax=4.2,
      axis lines=box,grid=major,grid style={gridgray!30},
      xlabel={$z$},ylabel={$\ell$},legend style={draw=none,font=\scriptsize}]
      \addplot[coursescore!88!ink,thick,domain=-4:4,samples=100]
        {ln(1+exp(-x))};
      \addlegendentry{\GryffondorName{} observé}
      \addplot[warning!88!ink,thick,domain=-4:4,samples=100]
        {ln(1+exp(x))};
      \addlegendentry{\SerpentardName{} observé}
    \end{axis}
  \end{tikzpicture}
  \[
    \ell(z,y)=\operatorname{softplus}(z)-yz.
  \]
\end{frame}

\begin{frame}[t]{Influence des profils en recouvrement}
  Avec les coefficients d'essai, le signe du score concorde avec la maison
  observée pour \MicroTrialAgree{} élèves et s'y oppose pour
  \MicroTrialDisagree{} :
  \[
    J=\frac1{\MicroNTrain}\sum_{i=1}^{\MicroNTrain}\ell(z_i,y_i)
      =\MicroTrialLoss.
  \]
  \vspace{3mm}
  \begin{columns}[T,onlytextwidth]
    \column{.47\textwidth}
      \begin{coursecallout}{rulecolor!88!ink}{Signe concordant}
        \MicroTrialBestName{} : \(z=\MicroTrialBestScore\), \(y=1\), perte
        \(\MicroTrialBestLoss\).
      \end{coursecallout}
    \column{.47\textwidth}
      \begin{coursecallout}{warning!88!ink}{Signe discordant}
        \MicroTrialWorstName{} : \(z=\MicroTrialWorstScore\), \(y=1\), perte
        \(\MicroTrialWorstLoss\), soit un rapport de \MicroTrialLossRatio.
      \end{coursecallout}
  \end{columns}
  \vfill
  Les observations situées dans l'agrégat de la maison opposée interdisent
  toute séparation linéaire parfaite et concentrent l'essentiel de la perte.
\end{frame}

% ---------------------------------------------------------------------------
\coursepart{Calculer le gradient}{Gradient}{coursegradient}

\begin{frame}[t]{Le gradient est une moyenne de résidus pondérés}
  \vspace{4mm}
  \[
    \frac{\partial\ell_i}{\partial z_i}=p_i-y_i.
  \]
  Comme \(z_i=w_0+w_{\mathrm{pot}}x_{i,\mathrm{pot}}
                   +w_{\mathrm{vol}}x_{i,\mathrm{vol}}\),
  \[
  \boxed{
    \nabla J(w)=\frac1n\sum_{i=1}^{n}(p_i-y_i)
      \begin{pmatrix}1\\x_{i,\mathrm{pot}}\\x_{i,\mathrm{vol}}\end{pmatrix}
    =\frac1nX^\top(p-y).}
  \]
\end{frame}

\begin{frame}[t]{Gradient au point d'initialisation}
  \[
    w^{(0)}=(0,0,0)\quad\Longrightarrow\quad z_i=0,
    \quad p_i=\frac12.
  \]
  Les \MicroNTrain{} observations donnent exactement
  \[
    \sum_i(p_i-y_i)=\MicroGradientZeroSumOne,\qquad
    \sum_i(p_i-y_i)x_{i,\mathrm{pot}}=\MicroGradientZeroSumPotions,\qquad
    \sum_i(p_i-y_i)x_{i,\mathrm{vol}}=\MicroGradientZeroSumVol.
  \]
  Donc
  \[
    \boxed{\nabla J(w^{(0)})=\MicroGradientZero.}
  \]
  La première composante est non nulle : les maisons ne sont pas également
  représentées, l'ordonnée à l'origine bouge dès le premier pas.
\end{frame}

\begin{frame}[t]{Contribution de cinq observations au gradient}
  \centering
  \renewcommand{\arraystretch}{1.22}
  \begin{tabular}{lrrrrr}
    \toprule
    élève & $y$ & $p-y$ à $w=0$ & $x_{\mathrm{pot}}$ &
      $(p-y)x_{\mathrm{pot}}$ & $(p-y)x_{\mathrm{vol}}$\\
    \midrule
    \MicroResidualRows
    \bottomrule
  \end{tabular}
  \vfill
  Les contributions sont additionnées avant la division par
  \(n=\MicroNTrain\).
\end{frame}

\begin{frame}[t]{Calcul vectoriel du gradient}
  \vspace{3mm}
  \[
    \begin{aligned}
      z&=Xw,\\
      p&=\sigma(z),\\
      r&=p-y,\\
      \nabla J(w)&=\frac1{\MicroNTrain}X^\top r.
    \end{aligned}
  \]
  \vfill
  \begin{coursecallout}{coursegradient!88!ink}{Dimensions}
    \(X\in\mathbb R^{\MicroNTrain\times3}\), \(w\in\mathbb R^3\),
    \(p,y,r\in\mathbb R^{\MicroNTrain}\) et
    \(\nabla J(w)\in\mathbb R^3\).
  \end{coursecallout}
\end{frame}

\begin{frame}[t]{Vérification du gradient par différence finie}
  \[
    \frac{\partial J}{\partial w_j}\approx
      \frac{J(w+h e_j)-J(w-h e_j)}{2h},\qquad h=10^{-5}.
  \]
  \centering
  \scriptsize
  \renewcommand{\arraystretch}{1.2}
  \begin{tabular}{crrr}
    \toprule
    composante $j$ & gradient analytique & différence finie & écart\\
    \midrule
    \MicroGradientCheckRows
    \bottomrule
  \end{tabular}
  \vfill
  Écart maximal : \(\MicroGradientMaxError\).
\end{frame}

% ---------------------------------------------------------------------------
\coursepart{Estimer les coefficients}{Mise à jour}{courseupdate}

\begin{frame}[t]{Première mise à jour des coefficients}
  \[
    w^{(t+1)}=w^{(t)}-\alpha\nabla J(w^{(t)}),\qquad \alpha=1.
  \]
  À partir de \(w^{(0)}=(0,0,0)\) :
  \[
    \begin{aligned}
      w^{(1)}&=(0,0,0)-\MicroGradientZero\\
             &=\boxed{\MicroFirstStep}.
    \end{aligned}
  \]
  \vspace{1mm}
  La perte moyenne passe de \(\MicroLossAtZero\) à \(\MicroLossAfterStep\) :
  une itération réduit le coût sans encore fixer de règle de décision.
\end{frame}

\begin{frame}[t]{Convergence de la descente de gradient}
  \centering
  \scriptsize
  \renewcommand{\arraystretch}{1.08}
  \begin{tabular}{rrrrrr}
    \toprule
    itération & $J$ & $w_0$ & $w_{\mathrm{pot}}$ & $w_{\mathrm{vol}}$ &
      $\lVert\nabla J\rVert$\\
    \midrule
    \MicroGDTableRows
    \bottomrule
  \end{tabular}
  \vfill
  Le critère \(\lVert\nabla J\rVert\leq10^{-12}\) est atteint à l'itération
  \(\MicroGDIterations\).
\end{frame}

\begin{frame}[t]{Ce que caractérise l'optimum}
  L'optimum n'admet pas de forme close : il est caractérisé par l'annulation du
  gradient. Les résidus \(p_i-y_i\) y sont de somme nulle et orthogonaux aux
  deux notes.
  \par
  \centering
  \small
  \renewcommand{\arraystretch}{1.2}
  \begin{tabular}{ccc}
    \toprule
    colonne de $X$ & condition & valeur atteinte\\
    \midrule
    \MicroFirstOrderRows
    \bottomrule
  \end{tabular}
  \normalsize
  \vfill
  \begin{coursecallout}{courseupdate!88!ink}{Unicité de l'optimum}
    Newton–Raphson converge en une dizaine d'itérations, la descente de
    gradient à pas constant en \MicroGDIterations{}. Les deux méthodes
    atteignent le même vecteur : la perte logistique est strictement convexe
    sous recouvrement des deux maisons.
  \end{coursecallout}
\end{frame}

\begin{frame}[t]{Paramètres standardisés et paramètres bruts}
  \begin{columns}[T,onlytextwidth]
    \column{.48\textwidth}
      \begin{coursecallout}{courseupdate!88!ink}{Variables standardisées}
        \begin{adjustbox}{max width=\linewidth,center}
          $(w_0,w_{\mathrm{pot}},w_{\mathrm{vol}})=\MicroWDecimal$
        \end{adjustbox}
        \par\smallskip
        {\scriptsize Le coefficient de \PotionsName{} vaut deux fois celui de
        \VolName{} : un écart-type de \PotionsName{} déplace le score deux
        fois plus qu'un écart-type de \VolName.}
      \end{coursecallout}
    \column{.48\textwidth}
      \begin{coursecallout}{coursescore!88!ink}{Notes brutes}
        \begin{adjustbox}{max width=\linewidth,center}
          $(b,\beta_{\mathrm{pot}},\beta_{\mathrm{vol}})=\MicroRawBetaDecimal$
        \end{adjustbox}
        \par\smallskip
        {\scriptsize Le rapport s'inverse sur les notes brutes : un point de
        \PotionsName{} équivaut à dix points de \VolName, les deux échelles n'étant pas
        commensurables.}
      \end{coursecallout}
  \end{columns}
  \vfill
  \[
    \boxed{\MicroRawScore}
  \]
  L'ordonnée à l'origine n'est pas nulle et les deux poids n'ont pas la même
  intensité : la frontière ne passe ni par le barycentre ni le long d'une
  bissectrice. \(J(w^\star)=\MicroTrainLoss\).
\end{frame}

\begin{frame}[t]{Le recouvrement rend l'optimum fini}
  \centering
  \begin{tikzpicture}
    \begin{axis}[width=.74\textwidth,height=4.8cm,xmin=0,xmax=2.6,
      ymin=.45,ymax=.72,axis lines=box,grid=major,grid style={gridgray!30},
      xlabel={$t$ dans $w=t\,w^\star$},ylabel={$J(t)$},
      tick label style={font=\scriptsize},label style={font=\scriptsize}]
      \addplot[courseupdate!88!ink,very thick,domain=0:2.6,samples=140]
        {\MicroLossCurveExpr};
      \addplot[only marks,mark=*,mark size=2.8pt,warning]
        coordinates {(1,\MicroTrainLossPlain)};
      \node[font=\scriptsize,anchor=north west,text=warning]
        at (axis cs:1,\MicroTrainLossPlain) {$t=1$};
    \end{axis}
  \end{tikzpicture}
  \begin{releve}
    \MicroWrongSideNames{} se situent du côté opposé à leur maison. En
    l'absence de ces \MicroWrongSideCount{} observations les deux agrégats
    seraient linéairement séparables : la perte décroîtrait vers zéro lorsque
    \(t\) croît, les coefficients divergeraient et l'optimum ne serait pas
    atteint.
  \end{releve}
\end{frame}

\begin{frame}[t]{Frontière du modèle ajusté}
  \begin{columns}[c,onlytextwidth]
    \column{.60\textwidth}
      \begin{tikzpicture}
        \begin{axis}[rawcloud,legend columns=3]
          \addplot[only marks,mark=square*,mark size=2.6pt,coursescore!88!ink]
            coordinates {\MicroTrainGryffondorCoords};
          \addlegendentry{\GryffondorName}
          \addplot[only marks,mark=*,mark size=2.6pt,warning!88!ink]
            coordinates {\MicroTrainSerpentardCoords};
          \addlegendentry{\SerpentardName}
          \microboundary{ink!82,very thick};
          \addlegendentry{$p=0{,}50$}
          \MicroWrongSideLabels
        \end{axis}
      \end{tikzpicture}
    \column{.36\textwidth}
      \begin{coursecallout}{warning!88!ink}{Observations mal classées}
        \MicroWrongSideNames.
      \end{coursecallout}
      \par\smallskip
      {\small \(\MicroBoundaryRaw\)\par\smallskip
      Le modèle ajusté affecte ces \MicroWrongSideCount{} observations à la
      maison opposée.}
  \end{columns}
\end{frame}

% ---------------------------------------------------------------------------
\coursepart{Définir la règle de décision}{Décision}{coursedecision}

\begin{frame}[t]{De la probabilité à la maison prédite}
  \vspace{4mm}
  \[
    \widehat y_i=\mathbf1[p_i\geq\tau].
  \]
  \begin{columns}[T,onlytextwidth]
    \column{.47\textwidth}
      \begin{coursecallout}{warning!88!ink}{$p_i<\tau$}
        Maison prédite : \SerpentardName.
      \end{coursecallout}
    \column{.47\textwidth}
      \begin{coursecallout}{coursescore!88!ink}{$p_i\geq\tau$}
        Maison prédite : \GryffondorName.
      \end{coursecallout}
  \end{columns}
  \vfill
  Le seuil agit après l'estimation : le modifier ne recalcule ni la
  standardisation, ni les coefficients.
\end{frame}

\begin{frame}[t]{Distribution de la perte sur le groupe d'estimation}
  \[
    \ell_i=-\ln p_i \text{ si \GryffondorName},
    \qquad
    \ell_i=-\ln(1-p_i) \text{ si \SerpentardName}.
  \]
  \vspace{2mm}
  \begin{columns}[T,onlytextwidth]
    \column{.31\textwidth}
      \begin{coursecallout}{rulecolor!88!ink}{Au centre de son agrégat}
        \MicroCheapName{} : \(p=\MicroCheapProbability\), perte
        \(\MicroCheapLoss\).
      \end{coursecallout}
    \column{.31\textwidth}
      \begin{coursecallout}{coursedecision!88!ink}{Au voisinage de la frontière}
        \MicroBorderName{} : \(p=\MicroBorderProbability\), perte
        \(\MicroBorderLoss\).
      \end{coursecallout}
    \column{.31\textwidth}
      \begin{coursecallout}{warning!88!ink}{Dans l'agrégat opposé}
        \MicroExpensiveName{} : \(p=\MicroExpensiveProbability\), perte
        \(\MicroExpensiveLoss\).
      \end{coursecallout}
  \end{columns}
  \vfill
  \centering
  La perte moyenne \(J(w^\star)=\MicroTrainLoss\) est concentrée sur un petit
  nombre d'observations : elle mesure le coût du recouvrement.
\end{frame}

\begin{frame}[t]{Résultats sur le groupe d'estimation}
  \centering
  \scriptsize
  \renewcommand{\arraystretch}{.96}
  \begin{tabular}{cllrrrr}
    \toprule
    ID & élève & maison obs. & score & prob. & préd. 0,50 & préd. 0,65\\
    \midrule
    \MicroTrainCalcRows
    \bottomrule
  \end{tabular}
\end{frame}

\begin{frame}[t]{Matrice de confusion sur le groupe d'estimation}
  \centering
  \begin{tikzpicture}[x=1.45cm,y=1.02cm,
    cell/.style={draw=white,line width=1pt,minimum width=2.8cm,
      minimum height=1.35cm,font=\Large\bfseries},
    lab/.style={font=\small\bfseries,text=ink!75}]
    \node[lab,align=center] at (1,3.35) {\SerpentardName\\prédit};
    \node[lab,align=center] at (3,3.35) {\GryffondorName\\prédit};
    \node[lab,align=right,anchor=east] at (-.05,2.25) {\SerpentardName\\observé};
    \node[lab,align=right,anchor=east] at (-.05,.75) {\GryffondorName\\observé};
    \node[cell,fill=rulecolor!18,text=rulecolor!82!ink] at (1,2.25)
      {VN $=\MicroTrainTN$};
    \node[cell,fill=warning!18,text=warning!88!ink] at (3,2.25)
      {FP $=\MicroTrainFP$};
    \node[cell,fill=warning!18,text=warning!88!ink] at (1,.75)
      {FN $=\MicroTrainFN$};
    \node[cell,fill=rulecolor!18,text=rulecolor!82!ink] at (3,.75)
      {VP $=\MicroTrainTP$};
  \end{tikzpicture}
  \vfill
  \[
    \text{exactitude}
      =\frac{\MicroTrainTP+\MicroTrainTN}{\MicroNTrain}=\MicroTrainAccuracy,
    \qquad
    \text{prédire toujours \GryffondorName}
      =\frac{\MicroTrainGryffondor}{\MicroNTrain}=\MicroTrainBaseline.
  \]
\end{frame}

\begin{frame}[t]{Le seuil déplace la frontière sans réestimer le modèle}
  \begin{columns}[c,onlytextwidth]
    \column{.60\textwidth}
      \begin{tikzpicture}
        \begin{axis}[rawcloud,legend columns=2]
          \addplot[only marks,mark=square*,mark size=2.2pt,coursescore,forget plot]
            coordinates {\MicroTrainGryffondorCoords};
          \addplot[only marks,mark=*,mark size=2.2pt,warning,forget plot]
            coordinates {\MicroTrainSerpentardCoords};
          \microboundary{ink!80,thick};
          \addlegendentry{$\tau=0{,}50$}
          \microboundaryhigh{coursedecision!88!ink,dashed,thick};
          \addlegendentry{$\tau=0{,}65$}
        \end{axis}
      \end{tikzpicture}
    \column{.36\textwidth}
      \begin{coursecallout}{coursedecision!88!ink}{Translation, pas rotation}
        \begin{adjustbox}{max width=\linewidth,center}
          $b+\beta_{\mathrm{pot}}P+\beta_{\mathrm{vol}}V
            \geq\operatorname{logit}(\tau)$
        \end{adjustbox}
      \end{coursecallout}
      \par\smallskip
      {\small Pente inchangée, ordonnée à l'origine relevée de
      \MicroBoundaryShift{} points de \VolName{} : relever le seuil à
      \(0{,}65\) revient à exiger un \VolName{} supérieur à niveau de
      \PotionsName{} constant.}
  \end{columns}
\end{frame}

% ---------------------------------------------------------------------------
\coursepart{Évaluer le modèle hors échantillon}{Évaluation}{courseevaluation}

\begin{frame}[t]{Les paramètres sont figés avant l'évaluation}
  \vspace{5mm}
  \centering
  \begin{adjustbox}{max width=\textwidth,center}
  \begin{tikzpicture}[node distance=5mm,
    box/.style={draw=ink!25,rounded corners=2pt,minimum height=16mm,
      align=center,font=\scriptsize,inner xsep=3mm},
    arr/.style={-{Latex[length=2mm]},thick,ink!60}]
    \node[box,fill=courseevaluation!14] (row) {une observation\\à évaluer};
    \node[box,fill=coursedata!15,right=of row] (impute)
      {$\widetilde P=\MicroMedianPotions$ ; $\widetilde V=\MicroMedianVol$\\si une note manque};
    \node[box,fill=courseprobability!15,right=of impute] (scale)
      {$\mu=(\MicroMeanPotions,\MicroMeanVol)$\\$s=(\MicroSdPotions,\MicroSdVol)$};
    \node[box,fill=courseupdate!15,right=of scale] (model)
      {$\MicroRawScore$\\coefficients figés};
    \node[box,fill=courseevaluation!20,right=of model] (out)
      {$p$, maison prédite, perte};
    \draw[arr] (row) -- (impute);
    \draw[arr] (impute) -- (scale);
    \draw[arr] (scale) -- (model);
    \draw[arr] (model) -- (out);
  \end{tikzpicture}
  \end{adjustbox}
  \vfill
  Aucune statistique et aucun coefficient ne sont réestimés avec les huit élèves.
\end{frame}

\begin{frame}[t]{Prédictions sur le groupe d'évaluation}
  \centering
  \scriptsize
  \renewcommand{\arraystretch}{1.11}
  \begin{tabular}{cllrrrr}
    \toprule
    ID & élève & maison obs. & score & prob. & préd. 0,50 & préd. 0,65\\
    \midrule
    \MicroTestCalcRows
    \bottomrule
  \end{tabular}
  \vfill
  Perte moyenne : \(J_{\mathrm{eval}}=\MicroTestLoss\).
\end{frame}

\begin{frame}[t]{Analyse locale autour du seuil de décision}
  \begin{columns}[c,onlytextwidth]
    \column{.60\textwidth}
      \begin{tikzpicture}
        \begin{axis}[rawcloud,legend columns=2]
          \addplot[only marks,mark=square,mark size=3pt,very thick,coursescore!88!ink]
            coordinates {\MicroTestGryffondorCoords};
          \addlegendentry{\GryffondorName}
          \addplot[only marks,mark=o,mark size=3pt,very thick,warning!88!ink]
            coordinates {\MicroTestSerpentardCoords};
          \addlegendentry{\SerpentardName}
          \microboundary{ink!78,very thick};
          \addlegendentry{$\tau=0{,}50$}
          \microboundaryhigh{coursedecision!88!ink,dashed,thick};
          \addlegendentry{$\tau=0{,}65$}
          \MicroThresholdBandLabels
        \end{axis}
      \end{tikzpicture}
    \column{.36\textwidth}
      \begin{coursecallout}{coursedecision!88!ink}{Entre les deux droites}
        Olivier : \(p=\MicroProbabilityOlivier\).\par
        Daphné : \(p=\MicroProbabilityDaphne\).
      \end{coursecallout}
      \par\smallskip
      {\small Seules ces deux observations changent de classe prédite lors du
      passage de \(\tau=0{,}50\) à \(\tau=0{,}65\).}
  \end{columns}
\end{frame}

\begin{frame}[t]{Effet du seuil sur la matrice de confusion}
  \begin{columns}[T,onlytextwidth]
    \column{.48\textwidth}
      \centering
      {\bfseries\color{coursedecision!88!ink} $\tau=0{,}50$}\par\medskip
      \begin{tabular}{c|cc}
        & Serp. prédit & Gryff. prédit\\\hline
        Serp. obs. & $\MicroTestTN$ & \cellcolor{warning!18}$\MicroTestFP$\\
        Gryff. obs. & \cellcolor{warning!18}$\MicroTestFN$ &
          \cellcolor{rulecolor!18}$\MicroTestTP$
      \end{tabular}
      \par\medskip
      Daphné est un faux positif, Alicia un faux négatif.
    \column{.48\textwidth}
      \centering
      {\bfseries\color{coursedecision!88!ink} $\tau=0{,}65$}\par\medskip
      \begin{tabular}{c|cc}
        & Serp. prédit & Gryff. prédit\\\hline
        Serp. obs. & \cellcolor{rulecolor!18}$\MicroTestTNHigh$ &
          $\MicroTestFPHigh$\\
        Gryff. obs. & \cellcolor{warning!18}$\MicroTestFNHigh$ &
          $\MicroTestTPHigh$
      \end{tabular}
      \par\medskip
      Plus aucun faux positif ; Olivier devient un faux négatif.
  \end{columns}
  \vfill
  L'exactitude reste \(\MicroTestCorrect/\MicroNTest\) ; la nature de l'erreur
  change. Le seuil élevé achète de la précision
  (\(\MicroTestPrecisionTauHigh\)) au prix du rappel
  (\(\MicroTestRecallTauHigh\)).
\end{frame}

\begin{frame}[t]{Comparaison des métriques selon le seuil}
  \centering
  \renewcommand{\arraystretch}{1.35}
  \begin{tabular}{rrrrrr}
    \toprule
    seuil & exactitude & précision & rappel & spécificité & F1\\
    \midrule
    \MicroTestSummaryRows
    \bottomrule
  \end{tabular}
  \vfill
  À \(\tau=0{,}50\) : précision \(=\MicroTestPrecision\),
  rappel \(=\MicroTestRecall\), spécificité \(=\MicroTestSpecificity\),
  \(F_1=\MicroTestFScore\).
\end{frame}

\begin{frame}[t]{Incertitude d'évaluation sur \MicroNTest{} observations}
  \vspace{3mm}
  \[
    \widehat{\mathrm{acc}}
      =\frac{\MicroTestCorrect}{\MicroNTest}=\MicroTestAccuracy.
  \]
  L'intervalle de Wilson à 95\% vaut
  \[
    \boxed{[\MicroTestWilsonLow\,;\,\MicroTestWilsonHigh]}.
  \]
  \begin{vigilance}[Portée de l'étude de cas]
    Les \MicroTotalStudents{} observations sont synthétiques et choisies pour
    rendre chaque transformation vérifiable à la main. Elles illustrent une
    méthode ; elles ne mesurent pas la performance attendue sur une population
    réelle.
  \end{vigilance}
\end{frame}
